% Part: computability
% Chapter: recursive-functions
% Section: sequences

\documentclass[../../../include/open-logic-section]{subfiles}

\begin{document}

\olfileid{cmp}{rec}{seq}
\olsection{لړۍ}

د بنسټيزو بازګشتي تابعو سټ د پام وړ ټينګ دے۔ خو کله چې د \emph{لړيو} د سمبالولو يوه مناسبه طريقه جوړه کړو، لا ډېر څه کولے شو۔ د طبيعي عددونو متناهي لړۍ به له طبيعي عددونو سره په لاندې ډول وپېژنو: لړۍ~$\langle a_0, a_1, a_2, \dots, a_k \rangle$ له دې عدد سره سمون لري:
\[
p_0^{a_0+1} \cdot p_1^{a_1+1} \cdot p_2^{a_2+1} \cdot \dots \cdot p_k^{a_k+1}.
\]
توانونو ته يو ځکه زياتوو چې مثلاً لړۍ~$\langle 2, 7, 3\rangle$ او $\langle 2, 7, 3, 0, 0 \rangle$ بېل عددي کوډونه ولري۔ $0$ او~$1$ دواړه د تشې لړۍ کوډونه نيولے شو؛ د کره‌والي دپاره دې~$\emptyseq$ د~$0$ ښودنه وکړي۔

د لړيو دا کوډول د حساب د تش په نامه بنسټيزې قضيې له امله کار کوي: هر طبيعي عدد~$n \ge 2$ په پوره يوه طريقه په لاندې بڼه ليکل کېدے شي:
\[
n = p_0^{a_0} \cdot p_1^{a_1} \cdot \dots \cdot p_k^{a_k}
\]
چې~$a_k \ge 1$ وي۔ دا ضمانت کوي چې نقشه~$\tuple{}(a_0, \dots, a_k) = \tuple{a_0, \dots, a_k}$ يو په يو ده: بېلې لړۍ بېلو عددونو ته لېږل کېږي، او له هر عدد سره تر ټولو زياته يوه لړۍ سمون لري۔

اوس به وښيو چې د يوې لړۍ د اوږدوالي موندل، د هغې د~$i$م غړي موندل، د لړۍ په پاے کښې د يوه غړي زياتول، او د دوو لړيو نښلول ټولې بنسټيزې بازګشتي چارې دي۔

\begin{prop}
  تابع~$\len{s}$، چې د لړۍ~$s$ اوږدوالی راګرځوي، بنسټيزه بازګشتي ده۔
\end{prop}

\begin{proof}
  اړيکه~$R(i, s)$ داسې تعريف کړئ:
  \[
  R(i, s) \text{ هله او يوازې هله چې } p_i \mid s \land p_{i+1} \nmid s.
  \]
  $R$ په څرګند ډول بنسټيزه بازګشتي ده۔ کله چې~$s$ د يوې ناتشې لړۍ کوډ وي، يعنې
  \[
  s = p_0^{a_0+1} \cdot \dots \cdot p_{k}^{a_{k}+1},
  \]
  $R(i,s)$ هله صدق کوي چې~$p_i$ تر ټولو لوے اوليه عدد وي چې $p_i \mid s$، يعنې~$i = k$ وي۔ نو د~$s$ اوږدوالی هله او يوازې هله~$i+1$ دے چې~$p_i$ تر ټولو لوے اوليه عدد وي چې~$s$ وېشي؛ ځکه دا تعريف ورکولے شو:
  \[
  \len{s} =
  \begin{cases}
    0 & \text{که $s = 0$ يا $s = 1$ وي} \\
    1 + \bmin{i < s}{R(i, s)} & \text{که نۀ}
  \end{cases}
  \]
  دلته محدوده کمينه‌موندنه کارولے شو، ځکه کله چې~$s$ د يوې لړۍ کوډ وي، يوازې يو~$i$ اړيکه~$R(i,s)$ پوره کوي، او که~$i$ موجود وي نو پخپله له~$s$ څخه کوچنے دے۔
\end{proof}

\begin{prop}
  تابع~$\fn{append}(s,a)$، چې د لړۍ~$s$ په پاے کښې د~$a$ د ورزياتولو پايله راګرځوي، بنسټيزه بازګشتي ده۔
\end{prop}

\begin{proof}
  $\fn{append}$ داسې تعريفېدے شي:
  \[
  \fn{append}(s,a) =
  \begin{cases}
    2^{a+1} & \text{که $s = 0$ يا $s = 1$ وي} \\
    s \cdot p_{\len{s}}^{a+1} & \text{که نۀ۔}
  \end{cases}
  \]
\end{proof}

\begin{prop}
  تابع~$\fn{element}(s,i)$، چې د~$s$ $i$م غړے راګرځوي (لومړني غړي ته~$0$م وايي)، او که~$i$ د~$s$ له اوږدوالي څخه لوے يا ورسره برابر وي نو~$0$ راګرځوي، بنسټيزه بازګشتي ده۔
\end{prop}

\begin{proof}
پام وکړئ چې~$a$ هله او يوازې هله د~$s$ $i$م غړے دے چې $p_i^{a+1}$ د~$p_i$ تر ټولو لوے توان وي چې~$s$ وېشي؛ يعنې $p_i^{a+1} \mid s$ وي خو~$p_i^{a+2} \nmid s$ وي۔ نو:
  \[
  \fn{element}(s,i) =
  \begin{cases}
    0 & \mbox{که $i \geq \len{s}$ وي} \\
    \bmin{a < s}{(p_i^{a+2} \nmid s)} & \text{که نۀ۔}
  \end{cases}
  \]
\end{proof}

د پورته تعريف شوو تابعو د رسمي نومونو پر ځاے لا لنډې نښې کاروو۔ د~$\fn{element}(s,i)$ پر ځاے به~$(s)_i$ کاروو، او $\tuple{s_0, \dots, s_k}$ به د لاندې عبارت لنډون وي:
\[
\fn{append}(\fn{append}(\dots \fn{append}(\emptyseq,s_0) \dots),s_k).
\]
پام وکړئ که~$s$ د~$k$ اوږدوالی ولري، نو د~$s$ غړي $(s)_0$، \dots،~$(s)_{k-1}$ دي۔

\begin{prop}
تابع~$\fn{concat}(s,t)$، چې دوه لړۍ سره نښلوي، بنسټيزه بازګشتي ده۔
\end{prop}

\begin{proof}
  داسې تابع~$\fn{concat}$ غواړو چې دا خاصيت ولري:
  \[
    \fn{concat}(\tuple{a_0, \dots, a_k}, \tuple{b_0, \dots, b_l}) = \tuple{a_0, \dots, a_k, b_0, \dots, b_l}.
  \]
  مرستندويه تابع~$\fn{hconcat}(s,t,n)$ کاروو چې د~$t$ لومړۍ~$n$ نښې له~$s$ سره نښلوي۔ دا تابع په بنسټيز بازګښت داسې تعريفېدے شي:
  \begin{align*}
    \fn{hconcat}(s,t,0) & = s\\
    \fn{hconcat}(s,t,n+1) & = \fn{append}(\fn{hconcat}(s,t,n),(t)_n)
    \intertext{بيا~$\fn{concat}$ داسې تعريفولے شو:}
    \fn{concat}(s,t) & = \fn{hconcat}(s,t,\len{t}).
  \end{align*}
\end{proof}

د~$\fn{concat}(s,t)$ پر ځاے به~$s \concat t$ ليکو۔

دا به راته ګټوره وي چې د يوې لړۍ عددي کوډ د هغې د اوږدوالي او تر ټولو لوے غړي له مخې محدود کړے شو۔ فرض کړئ~$s$ د~$k$ اوږدوالي يوه لړۍ ده او هر غړے ئې له کوم عدد~$x$ څخه کوچنے يا ورسره برابر دے۔ که لړۍ ناتشه وي، نو~$s$ تر ټولو زيات~$k$ اوليه عوامل لري، هر يو ئې تر ټولو زيات~$p_{k-1}$ دے، او هر يو د~$s$ په اوليه تجزيه کښې تر ټولو زيات د~$x+1$ توان ته پورته شوے دے۔ نو داسې تعريف کوو:
\[
\fn{sequenceBound}(x,k) =
\fn{cond}(k,1,p_{\fn{pred}(k)}^{k \cdot (x+1)})
\]
بيا د پورته بيان شوې لړۍ~$s$ عددي کوډ تر ټولو زيات $\fn{sequenceBound}(x,k)$ دے۔
\textbf{د ژباړې د نسخې سمون (OLCMP-011):} سرچينې د لړۍ په حد کښې د تشې لړۍ دپاره د منفي شاخص اوليه عدد کاراوه؛ ورپسې پلټنې هم نوماند له حد څخه سخت کوچنے غوښته او له حد سره برابر قانوني کوډ ئې پرېښود۔ دلته شرطي تابع د تشې لړۍ حد يو ټاکي او پلټنه تر حد پورې ټول‌شموله شوې ده۔

پر لړيو د داسې حد په لرلو نوې تابعې په محدودې پلټنې تعريفولے شو۔ مثلاً~$\fn{concat}$ په محدودې پلټنې تعريفولے شو۔ يوازې د هغه څيز (د نښلول شوې لړۍ د عدد) يو بنسټيز بازګشتي \emph{ځانګړتيايي شرط} ليکل او د پلټنې حد ټاکل پکار دي۔ لاندې تعريف کار کوي:
\begin{align*}
  \fn{concat}(s,t) = {} & \bmin{v \leq \fn{sequenceBound}(s+t,\len{s} + \len{t})}{} \\
  & \quad(\len{v} = \len{s} + \len{t} \land {}\\
  & \qquad \bforall{i < \len{s}}{((v)_i = (s)_i)} \land {} \\
  & \qquad \bforall{j < \len{t}}{((v)_{\len{s}+j} = (t)_j)})
\end{align*}
\textbf{د ژباړې د نسخې سمون (OLCMP-012):} سرچينې د دويمې لړۍ د غړو شرط د لومړۍ لړۍ پر غړو د کلي کميت ټاکونکي دننه ايښی و؛ نو د تشې لومړۍ لړۍ په حالت کښې هغه شرط تش‌صدقه کېده او د دويمې لړۍ هېڅ غړے ئې نۀ ټاکه۔ دلته د دواړو لړيو شرطونه د اوږدوالي له شرط سره په خپلواکو عطفونو تړل شوي دي۔

\begin{prob}
وښيئ چې داسې بنسټيزه بازګشتي تابع~$\fn{sconcat}(s)$ شته چې دا خاصيت ولري:
\[
\fn{sconcat}(\tuple{s_0, \dots, s_k}) = s_0 \concat \dots \concat s_k.
\]
\end{prob}

\begin{prob}
وښيئ چې داسې بنسټيزه بازګشتي تابع~$\fn{tail}(s)$ شته چې دا خاصيت ولري:
\begin{align*}
  \fn{tail}(\emptyseq) & = 0 \text{ او}\\
  \fn{tail}(\tuple{s_0, \dots, s_{k}}) & = \tuple{s_1, \dots, s_{k}}.
\end{align*}
\end{prob}

\begin{prop}
  \ollabel{prop:subseq}
  تابع~$\fn{subseq}(s, i, n)$ چې د~$s$ له~$i$م غړي څخه پيلېدونکې د $n$ اوږدوالي فرعي لړۍ راګرځوي، بنسټيزه بازګشتي ده۔
\end{prop}

\begin{proof}
  تمرين۔
\end{proof}

\begin{prob}
  \olref[cmp][rec][seq]{prop:subseq} ثابت کړئ۔
\end{prob}

\end{document}
